\documentclass[../main.tex]{subfiles} % Due to use of package subfiles

\pagestyle{empty} % No page numbers on page

%%%%%%%%%%%%%%%%%%%%%%%%%%%%%%%%%%%%%%%%%%%%%%%%%%%%%%%%%%%%%%%%%%%%%%%%%%%%%%%%

\begin{document}

\clearpage % Needs to be here if bookmarks shall be used
\pdfbookmark[0]{Back cover}{back cover} % Bookmark in PDF

\newgeometry{lmargin=4.5cm, rmargin=4.5cm, tmargin=3.8cm, bmargin=2.8cm} % Sets the margins for the front page

% Logo
\begin{figure}[ht!] % Begin front page image
    \begin{center}
        \includegraphics[width=3.65cm]{afterText/backCoverImages/auseglBlue.eps}
    \end{center}
\end{figure} % End front page image

\vspace{1em}
\vfill

Quantum computing and its promise of significant speedup for certain problems
is making the interest and research in the field ever growing. Superconducting
circuits offer a promising platform for the implementation of computing,
as they are versatile and potent with respect to the quantum effects and
parameter ranges they offer. The purpose of quantum computers as envisaged
by Feynman, however, is to perform analogue simulations of physics for the
sake of research. The idea is to construct a quantum system in the laboratory
which behaves like some interesting system in Nature. The lab-system
might then be more accessible and easier to control such that consistent and
systematic research can take place.

Something especially suited for analogue simulation is spin models, as spins
and qubits (the fundamental unit of a quantum computer) are formally identical.
Many theories can in turn be mapped on to spin models, or at least spin
models with qualitatively similar properties can be constructed. Among these
are lattice gauge theories, the discretized version of gauge invariant quantum
field theories. Many fundamental results have been found by considering especially
quantum chromodynamics, the gauge theory which describes the strong
nuclear force as a lattice gauge theory. One thing that is not yet understood,
however, is the so-called strong CP problem. It is theoretically allowed for the
strong force to break CP-symmetry, however the vacuum angle, a parameter
which quantifies this breaking, is experimentally found to be extremely small,
for which there is not yet an accepted theoretical explanation.

for which there is not yet an accepted theoretical explanation. for which there is not yet an accepted theoretical explanation.

\vfill
\vspace{2em}

% Logo
\begin{figure}[ht!] % Begin front page image
    \begin{center}
        \includegraphics[width=\textwidth]{afterText/backCoverImages/AUDepartmentLogoBlue.eps}
    \end{center}
\end{figure} % End front page image

\restoregeometry % Sets the geometry (margins) back to the in preamble defined

\end{document}

%%%%%%%%%%%%%%%%%%%%%%%%%%%%%%%%%%%%%%%%%%%%%%%%%%%%%%%%%%%%%%%%%%%%%%%%%%%%%%%%